\documentclass[11pt]{article}

\usepackage{amsmath, amssymb, amsthm}
\usepackage{bm}
\usepackage{geometry}
\usepackage{hyperref}
\usepackage{natbib}
\usepackage{xcolor}
\usepackage{graphicx}

\geometry{margin=1in}

\newcommand{\given}{\,|\,}
\newcommand{\Normal}{\mathcal{N}}
\newcommand{\Lognormal}{\text{Lognormal}}
\newcommand{\HalfNormal}{\text{HalfNormal}}
\newcommand{\Exponential}{\text{Exponential}}
\newcommand{\HalfCauchy}{\text{HalfCauchy}}
\newcommand{\Poisson}{\text{Poisson}}
\newcommand{\NegBin}{\text{NegBin}}
\newcommand{\Beta}{\text{Beta}}

\setlength\parindent{0pt}

\graphicspath{{fig/}}


\title{\texttt{SCARCHhierarSIR}\\Training and Forecasting Procedure}
\author{Tijs~W.~Alleman~\& Ana~I.~Bento}
\date{\today}

\begin{document}
\maketitle

% Introduction
% What is the hierarSIR model?

%-----------------------------------------------------------
\section{Overview}
%-----------------------------------------------------------

Our methodology consists of two distinct phases. \textbf{Training:} Embeds a forward-simulating SIR model in a Bayesian hierarchical inference framework. Season- and state-specific SIR parameters are modeled as exchangeable draws from shared hyperdistributions, whose posterior is inferred from historical data. \textbf{Forecasting:} Uses the inferred hyperdistributions as informative priors for a Bayesian update of the within-season parameters during an ongoing season, producing four-week-ahead predictive distributions, as well as estimates of peak magnitude and timing that are updated weekly as new data arrive.

%-----------------------------------------------------------
\section{Framework components}
\subsection{Foward-simulating SIR model}
%-----------------------------------------------------------

% TODO: add the subscripts i,j

\paragraph{Compartmental structure} The transmission of influenza in a given U.S.\ state or territory is governed by the following temporally-forced deterministic variant of the canonical SIR model \citep{keeling2008},
\begin{align}
    \lambda_{i,j}(t) &= \beta_0 [1 + \textcolor{red}{\Delta \beta_{i,j}(t)}]  I_{i,j}(t), \nonumber \\
    \dot{S}_{i,j}(t) &= - \lambda_{i,j}(t) S_{i,j}(t), \nonumber \\ 
    \dot{I}_{i,j}(t) &= \lambda_{i,j}(t) S_{i,j}(t) - \gamma I_{i,j}(t), \nonumber \\
    \dot{R}_{i,j}(t) &= \gamma I_{i,j}(t),
    \label{eq:ODE_simulation}
\end{align}
where $S_{i,j}(t) + I_{i,j}(t) + R_{i,j}(t) = 1$. In this model, $\lambda_{i,j}(t)$ is the force of infection at day $t$, in season $i$ and U.S.\ state or territory $j$. $\beta_0$ represents the baseline transmission coefficient and $\gamma=1/3.5~d^{\text{-}1}$ is the average duration of infectiousness, chosen to correspond to the CDC Yellow Book estimate of the most contagious phase of influenza infection \citep{CDCyellowbook}.\\

$\Delta \beta_{i,j}(t)$ represents a time-varying multiplicative modifier to the baseline transmission coefficient $\beta_0$, obtained by smoothing a piecewise-continuous trajectory of 32 week-long modifiers $\Delta \mathcal{B}_{i,j}(l)$ with a Gaussian filter ($\sigma = 1~\text{d.}$). For the influenza season of year $\{\text{X}-\text{X}+1\}$, the first modifier represents 1-7 October of year $\text{X}$ and the last modifier represents 1-7 May of year $\text{X}+1$. Outside of this range, there is no modification of the transmission coefficient and hence $\Delta \beta_{i,j}(t)=0$. This approach extends the flexibility of the SIR model, allowing it to better capture complex patterns in the observed incidence data, effectively acting as a discrepancy model in the transmission domain \citep{osthus2019}.

\paragraph{Observations} We model weekly new hospital admissions as partially observed realizations of infection incidence, which are treated as approximately contemporaneous with incidence. Let \( I_{i,j}^{\text{new}}(t) = \lambda_{i,j}(t) S_{i,j}(t) T \) denote the incidence of new infections at time \(t\), with $T$ equal to the total population. 
\begin{equation}\label{eq:observed_states}
    \text{H}^{\text{adm}}_{i,j}(t) 
    = \textcolor{red}{\rho_{i,j}} \int_{t-7}^{t} I^{\text{new}}_{i,j}(s - \theta)\, ds,
\end{equation}
where \( \rho_{i,j} \) is the ascertainment probabilities for hospital admissions.

\paragraph{Initial conditions} We model seasonal influenza on a season-by-season basis. For the influenza season of year $\{\text{X} \text{--} \text{X}+1\}$, simulations are started on October 1st of year $\text{X}$. We model the initial condition as follows,
\begin{align}\label{eq:initial_condition}
    S_{i,j}(0) &= 1- f_{I,i,j}-f_{R,i,j}, \nonumber \\
    I_{i,j}(0) &= \textcolor{red}{f_{I,i,j}}, \nonumber \\
    R_{i,j}(0) &= \textcolor{red}{f_{R,i,j}}. 
\end{align}
Here $f_{I,i,j}$ represents the initial fraction of the population of U.S.\ state or territory $j$ infected at the start of season $i$, and $f_{R,i,j}$ represents the initial fraction of the population with immunity to influenza.\\

Due to its simplicity, the model is integrated using the second-order accurate Heun's method \citep{heun1900}, implemented via the \texttt{diffrax} library in \texttt{jax} \citep{kidger2021}. To improve gradient stability, we employ a fixed step size of 3.5 days, which corresponds to the shortest timescale in the model. In practice, we omit simulating the recovered compartment to further accelerate computations.

%-----------------------------------------------------------
\subsection{Bayesian hierarchical inference framework}
%-----------------------------------------------------------

\subsubsection*{Foward-simulating SIR model parameters}
Three of the forward-simulating SIR parameters, highlighted in red in Eqs.~\eqref{eq:ODE_simulation},~\eqref{eq:observed_states},~\eqref{eq:initial_condition} carry a non-centered additive hierarchy across seasons and U.S\ states and territories: the case ascertainment rate $\rho$, the initial fractions of infected $f_I$ and recovered $f_R$ individuals in the population. The baseline transmission coefficient $\beta$ is fixed at $\beta_0 = 0.455$, corresponding to a basic reproduction number of $R_0 = 1.6$.\\

The case ascertainment rate in season $i$, U.S.\ state $j$ is modeled using the following additive structure with a log-link,
\begin{align}
	 \log \rho_{i,j} &= \log \bar{\rho} + \sigma^{(i)}_\rho\, u^{(i)}_{\rho} + \sigma^{(j)}_\rho\, u^{(j)}_{\rho}, \nonumber\\
    \log \bar{\rho} &\sim \Normal\!\left(\mu_{\rho},(1/2)^2\right), \nonumber\\
    \sigma^{(i)}_\rho,\, \sigma^{(j)}_\rho &\sim \HalfNormal\! \left( 1/5 \right), \nonumber\\
    u^{(i)}_{\rho}, u^{(j)}_{\rho} & \overset{\mathrm{iid}}{\sim} \Normal(0,1),
\end{align}
where the priors $\mu_{\rho} = \log \left( 0.5~\% \right)$ and $\sigma_{\rho} = 1/2$ center the hierarchy at 0.5~\% case ascertainment with 95~\% of the mass between $0.1~\%$ and $2.2~\%$. The halfnormal distribution on $\sigma_{\rho}^{(i)},~\sigma_{\rho}^{(j)}$ shrink season- and spatial effect sizes when the data does not support them.\\

The initial fraction of infected $f_I$ has small values ($f_I$ $\approx 10^{-5}$), and for this reason we use a log-link instead of the more natural logit-link function, because this results in better posterior geometry.\\

The initial fraction of recovered $f_R$ in season $i$, U.S.\ state $j$ is modeled using a logit-additive link to enforce the unit-interval constraint,
\begin{align}\label{eq:f_R}
	\text{logit}\,f_{R,i,j} &= \text{logit}\,\bar{f}_R + \sigma^{(i)}_{f_R}\, u^{(i)}_{f_R} + \sigma^{(j)}_{f_R}\, u^{(j)}_{f_R}, \nonumber\\
    \text{logit}\,\bar{f}_R &\sim \Normal\!\left(\text{logit}(0.4),\ 1\right), \nonumber\\
    \sigma^{(i)}_{f_R},\, \sigma^{(j)}_{f_R} &\sim \HalfNormal\!\left(1/5\right), \nonumber\\
    u^{(i)}_{f_R}, u^{(j)}_{f_R} & \overset{\mathrm{iid}}{\sim}\Normal(0,1),
\end{align}
The prior on $\text{logit}\,\bar{f}_R$ places the population-mean immunity at $40\pm10\%$, which results in an effective reproduction number of $R_e \approx 1$ at the start of the influenza season, anchoring the model, consistent with \texttt{hierarchSIR}.

%-----------------------------------------------------------
\subsubsection*{Spatially Correlated AR(1)--GARCH(1,1) Transmission Modifiers}\label{sec:argarch}
%-----------------------------------------------------------

\paragraph{Structure of the transmission modifiers} The effective transmission coefficient on day $t$ in season $i$ and state $j$ in the SIR model is,
\begin{equation}
\beta_{i,j}(t) = \beta_0 \left[ 1 + \Delta \beta_{i,j}(t) \right],
\end{equation}
where the modifier trajectory $\Delta \beta_{i,j}(t)$ is obtained by smoothing a piecewise-continuous trajectory of 32 week-long modifiers $\Delta \mathcal{B}_{i,j}(l)$ (indexed $l$) with a Gaussian filter ($\sigma = 1~\text{d.}$), which are in turn modeled as a stochastic structured time series,
\begin{equation}
    \Delta \mathcal{B}_{i,j} (l) = \Delta\bar{\mathcal{B}_{j}}(l) + z_{i,j}(l),
\end{equation}
where $\Delta\bar{\mathcal{B}_{j}}(l)$ is the seasonal mean modifier component in U.S.\ state or territory $j$ in week $l$, and $z_{i,j}(l)$ is the corresponding stochastic AR(1)--GARCH(1,1) component \citep{bollerslev1986}, which represents the current week's deviation from the seasonal mean modifier.


\paragraph{Seasonal mean modifiers}

To flexibly and continuously model temporal variation in the seasonal mean modifier while accounting for spatial dependence between states, we represent the state-specific modifier trajectories using a cubic B-spline basis and impose spatial correlation on the spline coefficients.\\


Specifically, let $\bm{B} \in \mathbb{R}^{L \times K}$ denote the cubic B-spline basis evaluated at the $L$ modifier weeks, with $K = 18$ basis functions and no intercept, and let $\bm{\zeta} \in \mathbb{R}^{K \times S}$ with $S$ the number of states denote the matrix of spline coefficients. For state $j$, the seasonal mean modifier at modifier week $l$ is given by,

\begin{equation}
\Delta \bar{\mathcal{B}}_j(l) = \sum_{k} B_{lk}\zeta_{kj},
\end{equation}


where $\zeta_{kj}$ denotes the coefficient of basis function $k$ for state $j$.\\

To induce spatial dependence among states, the spline coefficients $\bm{\zeta}$ are first represented as spatially correlated transformations of independent Laplace-distributed innovations. Specifically, for each spline basis function $k$, let,
\begin{equation}
u_{kj}
\overset{\mathrm{iid}}{\sim}
\text{Laplace}\left(0, (1/3)\right),
\end{equation}
and define the corresponding vector of spatially correlated spline coefficients as
\begin{equation}
\bm{\zeta}_k
=
\bm{L}(\psi_1)^{-\top}\bm{u}_k,
\end{equation}
where $\bm{\zeta}_k = (\zeta_{k1},\ldots,\zeta_{kS})^\top$ and $\bm{L}(\psi_1)$ is the Cholesky factor of the precision matrix,
\begin{equation}
\bm{Q}(\psi_1)
=
\bm{L}(\psi_1)\bm{L}(\psi_1)^\top.
\end{equation}
The precision matrix is defined as the weighted average between the spatially uncorrelated and correlated case, following \citep{leroux2000},
\begin{equation}
\bm{Q}(\psi_1)
=
(1-\psi_1)\bm{I}
+
\psi_1(\bm{D}-\bm{W}),
\qquad
\psi_1 \in (0,1),
\end{equation}
where $\psi_1 \sim \text{Beta}(3,3)$ controls the strength of spatial dependence, $\bm{W}$ denotes the binary adjacency matrix with $W_{jj'}=1$ if states $j$ and $j'$ share a border, and $\bm{D}$ denotes the corresponding diagonal degree matrix with entries $D_{jj}=\sum_{j'}W_{jj'}$.\\

This model is equivalent to a Leroux-style conditional autoregressive model \citep{leroux2000}, except that the underlying innovations $\bm{u}_k$ are assigned independent Laplace rather than Gaussian distributions. The Laplace distribution induces stronger shrinkage of small spline coefficients toward zero, leading to more stable estimates of the seasonal mean modifiers, while still permitting occasional larger deviations when supported by the data.


% https://www.youtube.com/watch?v=jzW40lOGebg&t=458s
% https://www.sciencedirect.com/science/article/pii/S1877584511000049
% https://link.springer.com/chapter/10.1007/978-1-4612-1284-3_4

\paragraph{Stochastic deviations}
The stochastic part $z_{i,j}(l)$ evolves according to an AR(1) process whose innovation variance follows a GARCH(1,1) model. For each season $i$, state $j$, and week $l$, the recursion is,
\begin{align}
	z_{i,j}(l) &= \phi \,z_{ij}(l-1) + \varepsilon_{ij}(l), \\ \label{eq:ar1}
	\varepsilon_{ij}(l) &= v_{ij}(l)\,\sigma_{ij}(l), \\
    \sigma^2_{ij}(l) &= \omega_G + \alpha_G \,\varepsilon^2_{ij}(l-1) + \beta_G\,\sigma^2_{ij}(l-1),
\end{align}
with initialization $z_{ij}(0) = 0$ and $\varepsilon_{ij}(0) = 0$. The innovation terms $\bm{v}_i(l) = (v_{i1}, \ldots, v_{iS})^T$ driving the system are spatially correlated across U.S.\ states and territories using a Leroux CAR model \citep{leroux2000},
\begin{equation}
    \bm{v}_i(l) = \bm{L}(\psi_2)^{-\top}\, \bm{u}_i(l), 
    \qquad u_{ij}(l) \overset{\mathrm{iid}}{\sim} \Normal(0, 1),
\end{equation}
resulting in $\bm{v}_i(l) \sim \Normal\big(\bm{0}, \bm{Q}(\psi_2)^{-1}\big)$.\\

Equation~\eqref{eq:ar1} represents the AR(1) dynamics governing deviations from the seasonal mean modifier, with persistence $\phi$. The GARCH(1,1) process governs the time-varying innovation variance, where $\omega_G$ controls the baseline variance, $\alpha_G$ captures the sensitivity of the system's volatility to new shocks, and $\beta_G$ controls persistence of its volatility.\\

We assume $\beta_G = 0$ for the sake of simplicity, and treat $\phi$, $\omega_G$ and $\alpha_G$ as hyperparameters that must be tuned to maximize 4-week ahead predictive accuracy by performing a grid search on validation data.


%-----------------------------------------------------------
\section{Training and Forecasting}
\subsection{Training}
%-----------------------------------------------------------
 
Collecting all within-season, within-state parameters into $\bm{\Theta}$ and all hyperparameters into $\bm{\Psi}$, the joint posterior sampled during training is,
\begin{equation}
    \mathcal{P}(\bm{\Theta}, \bm{\Psi} \given \bm{D}) \propto \mathcal{L}(\bm{D} \given \bm{\Theta})\, \mathcal{P}(\bm{\Theta} \given \bm{\Psi})\, \mathcal{P}(\bm{\Psi}),
\end{equation}
$\mathcal{P}(\bm{\Theta}, \bm{\Psi} \given \bm{D})$ is the posterior distribution of the parameters given the training data $\bm{D}$, $\mathcal{L}(\bm{D} \given \bm{\Theta}$) is a weighted Negative Binomial likelihood. $\mathcal{P}\left( \bm{\Theta} \given \bm{\Psi} \right)$ is the probability of observing the within-season parameters given the cross-season hyperparameters and captures the cross-season hierarchy of the model. $\mathcal{P} \left(\bm{\Psi} \right)$ are the prior distributions of the cross-season hyperparameters, which encode pre-existing knowledge about these parameters or enforce shrinkage.\\

A training dataset $\bm{D}$ spans at least one prior influenza season:
\begin{equation}
    \bm{D} = \bigl\{ D_{i,j}\ \forall\ i \in \bm{S},\ j \in \bm{T} \bigr\},
\end{equation}
Here, $\bm{S} = \{2023\text{-}2024,\, 2024\text{-}2025,\, 2025\text{-}2026\}$ is the set of training seasons, $\bm{T}$ is the set of U.S.\ states and territories, and $D_{i,j}$ is the NHSN HRD lab-confirmed hospital admissions for influenza in season $i$ and U.S.\ state or territory $j$. We represent an individual observed time series in season $i$ and U.,S\ state or territory $j$ as $y_{j,k}(t)$ and the corresponding simulation as $\hat{y}_{j,k}(t)$.\\
 
 
To match the forward simulation to the observed data, we employ a weighted Negative Binomial log-likelihood function for count data, 
\begin{equation}
    \log \mathcal{L}(\bm{D} \given \bm{\Theta}) = \sum_{i}\sum_{j}\sum_{t}\, w_{i,j} \cdot \log p \!\left(y_{i,j}(t) \given \mu = \hat{y}_{i,j}(t),\, \alpha^{-1}_{j}\right),
\end{equation}
where $p(y \given \mu, \alpha^{-1})$ is the negative binomial probability mass function with mean $\mu$ and overdispersion parameter $\alpha^{-1}$. We allow the overdispersion to vary by state,
\begin{equation}
    \alpha^{-1}_j \sim \HalfNormal\! \left( 0.003/3 \right),
\end{equation}
where the strong halfnormal prior shrinks the negative binomial to a Poisson distribution where possible, with the scale of $\sigma=0.001$ determined on a trial-and-error basis by assessing the coverage of the 95~\% confidence interval on the training data.\\

The season-state weights $w_{i,j}$ normalize all observed time series so that no single large-epidemic season or high-burden state dominates the likelihood,
\begin{equation}
    w_{i,j} = \frac{\left[ \text{mean}_{\{t\}} \big( \sqrt{y_{i,j}(t)}\, \big) \right]^{-1}}{\dfrac{1}{|\bm{S}| \cdot |\bm{T}|} \displaystyle\sum_{i'}\sum_{j'}  \left[ \text{mean}_{\{t\}} \left( \sqrt{y_{i',j'}(t)}\, \right) \right]^{-1}},
\end{equation}
alternatives to this normalization function use the average natural logarithm of the case counts or the maximum case count per season and state. Posterior sampling is performed using the No-U-Turn Sampler (NUTS) \citep{hoffman2011} implemented in \texttt{NumPyro} \citep{phan2019,bingham2019}.
 
%-----------------------------------------------------------
\subsection{Forecasting}
%-----------------------------------------------------------

During an ongoing season, a reduced model is run in which all global and state-level components, as well as the season-level hyperparameters are replaced by the posterior median found during the training phase. Hence, only the current season's offsets $\left(u^{(i)}_\rho, u^{(i)}_{f_I}, u^{(i)}_{f_R}, u^{(i)}_{\phi},\ \text{and}\ u^{(i)}_{\omega}\right)$ and the AR(1)--GARCH(1,1) stochastic trajectory $z_{i,j}(t)$ are sampled, conditioned on observations from the start of the season up to the forecast date.\\

Thus, the cross-season hyperdistributions $\bm{\Psi}$ are used as prior distributions for the inference of the forward simulation model's within-season parameters during the ongoing influenza season. The Bayesian approach to training and forecasting aligns naturally with the probabilistic nature of disease forecasting. Early in the season, when limited data are available, forecasts will be heavily informed by historical patterns, reflecting the average trends from previous years. As the season progresses and more data accumulate, the influence of the historical priors will naturally diminish. 

%-----------------------------------------------------------
\subsection{Hypertuning}
%-----------------------------------------------------------


The parameters governing the AR(1)--GARCH process, $\phi$, $\omega_G$, $\alpha_G$, and $\beta_G$, are treated as hyperparameters rather than inferred jointly with the remaining model parameters. This distinction is motivated by the weak identifiability of the GARCH parameters from the training likelihood. In particular, increasing the innovation variance allows the stochastic component to absorb increasingly large discrepancies between the forward simulation and the observations. Without external constraints on these parameters, the resulting model can attribute observation-level variation to process noise, leading to overfitting and limiting the interpretability and forecasting utility of the stochastic component. The temporal persistence parameter $\phi$ is more readily identified from the training data; however, its posterior estimates are relatively high (e.g., $\phi \approx 0.85$ when $\omega_G=0.005$), which may result in excessive persistence of stochastic deviations when extrapolated beyond the observation period.\\

Ideally, these hyperparameters would therefore be selected according to out-of-sample predictive performance rather than inferred from the training likelihood. We propose to perform leave-one-season-out cross-validation over the historical influenza seasons. For each held-out season, the model is trained on all remaining seasons, and forecasts are generated sequentially throughout the held-out season using only information available up to each forecast date. Predictive accuracy is then evaluated using the same scoring rules used for operational forecasting. This procedure is repeated over a predefined grid of values for $(\phi,\omega_G,\alpha_G,\beta_G)$, and the combination yielding the best average predictive performance across held-out seasons is selected.\\

For the 2026--2027 forecasting season, however, the available historical data are insufficient to perform this procedure reliably. The hierarchical model requires at least three seasons for training, leaving too few independent seasons for meaningful leave-one-season-out validation. Generating synthetic seasons to augment the validation set would introduce additional assumptions into the hyperparameter-selection procedure and could bias the resulting choices. We therefore fix the stochastic-process hyperparameters for the 2026--2027 season at $\alpha_G=\beta_G=0$, $\omega_G=0.005$, and $\phi=0.5$. These values provide a conservative stochastic component with moderate temporal persistence while avoiding the excessively large innovation variances and high persistence obtained when these parameters are estimated directly from the training likelihood.


%-----------------------------------------------------------
\section{Discussion of the changes to \texttt{hierarchSIR}}
%-----------------------------------------------------------

The \texttt{SCARCHhierarSIR} model extends the original \texttt{hierarchSIR} model in three important ways. First, whereas \texttt{hierarchSIR} treated influenza dynamics across U.S.\ states and territories as spatially independent, \texttt{SCARCHhierarSIR} incorporates spatial correlation within the Bayesian hierarchical inference layer. Importantly, spatial dependence is introduced only through the statistical model and does not alter the forward-simulating SIR model: the SIR dynamics remain independent across states and territories, with no explicit spatial transmission mechanism such as mobility between states.\\

Second, the computational framework has been substantially revised. Forward simulations are now performed using \texttt{diffrax}, a \texttt{JAX}-based library for differentiable numerical integration, enabling automatic differentiation, just-in-time compilation, and efficient batched simulation across states and seasons. The Bayesian inference layer is implemented in \texttt{NumPyro}, which is also built on \texttt{JAX}. This common computational framework allows the forward simulator and its gradients to be integrated directly into the probabilistic model without requiring an interface between separate programming environments. This in turn opens the door to efficient gradient-based exploration of the posterior using the No-U-Turn Sampler (NUTS) \citep{hoffman2011}. In contrast, the original \texttt{hierarchSIR} model implemented the forward simulation in \texttt{C++}, while the posterior model was implemented in \texttt{Python} and sampled using the Goodman--Weare ensemble sampler through \texttt{emcee}. The resulting \texttt{JAX}/\texttt{NumPyro} implementation substantially reduces the computational resources required for model training relative to \texttt{hierarchSIR}.\\

Third, and most importantly, \texttt{SCARCHhierarSIR} introduces a structured stochastic model for the time-varying transmission modifier. The modifier at each week is represented as the sum of a seasonal mean component and a stochastic deviation from that mean, with the latter generated by a spatially correlated AR(1)--GARCH(1,1) process. The AR(1) component captures the persistence of deviations from the typical seasonal trajectory, allowing an unusually high or low transmission period to influence subsequent weeks. The GARCH(1,1) component allows the magnitude of these stochastic deviations to vary over time, producing greater forecast uncertainty when the trajectory departs substantially from the seasonal mean. This provides a mechanism for the predictive distribution to adapt its width to the magnitude of observed departures from typical seasonal dynamics, which may improve probabilistic forecast performance as measured by the Weighted Interval Score (WIS). The contribution of this stochastic component to forecasting performance will be evaluated prospectively beginning with the 2027--2028 forecasting season. In contrast, the time-varying transmission modifier in \texttt{hierarchSIR} was modeled over a two-week window and subsequently reverted to the seasonal mean, without a persistent stochastic process governing deviations from that mean.

\clearpage
\bibliographystyle{plainnat}
\bibliography{bibliography}

\end{document}